\definecolor{c1}{HTML}{F2553B}
\definecolor{c2}{HTML}{5D6EEB}
\definecolor{c3}{HTML}{C9B6E9}
\definecolor{c4}{HTML}{39E0B0}
\definecolor{c5}{HTML}{F0A394}
\definecolor{c6}{HTML}{B8BFF1}
\definecolor{c7}{HTML}{F6C54D}
\definecolor{c8}{HTML}{F0A0F0}
\definecolor{c9}{HTML}{9EE06C}
\definecolor{c10}{HTML}{FF6FA9}
\definecolor{c11}{HTML}{22BDE6}
\definecolor{c12}{HTML}{F4A259}
\definecolor{c13}{HTML}{A375FF}
\definecolor{c14}{HTML}{2ED3A1}
\definecolor{c15}{HTML}{FF8AB3}
\definecolor{c16}{HTML}{C96A17}
\definecolor{c17}{HTML}{5E8E3E}
\definecolor{c18}{HTML}{A9A9A9}
\definecolor{c19}{HTML}{D91E36}
\begin{tikzpicture}

    % ======================
    % 参数设置
    % ======================
    \def\radius{4}        % 饼图半径
    \def\pctdistance{0.7} % 百分比文字距离圆心的比例 (对应 pctdistance=0.7)
    \def\currentangle{90} % 起始角度，12点钟方向 (对应 startangle=90)

    % ======================
    % 核心：绘制扇区和斜体文字
    % ======================
    \foreach \c/\l/\v in {
        c1/人与社会/18.61, c2/艺术与娱乐/12.08, c3/新闻/10.44, 
        c4/科学/8.79, c5/旅游与交通/6.43, c6/宠物与动物/5.79, 
        c7/健康/5.31, c8/敏感话题/4.91, c9/体育/4.29, 
        c10/书籍与文学/4.08, c11/法律与政府/3.86, c12/工作与教育/3.33, 
        c13/计算机与电子/2.01, c14/食品与饮料/1.96, c15/商业与工业/1.93, 
        c16/汽车与车辆/1.27, c17/游戏/1.14, c18/爱好与休闲/1.04, 
        c19/其他/2.74
    } {
        % 计算当前扇区的扫过角度、结束角度、以及中心角度
        \pgfmathsetmacro{\deltaangle}{\v * 360 / 100}
        \pgfmathsetmacro{\nextangle}{\currentangle - \deltaangle} % 顺时针递减
        \pgfmathsetmacro{\midangle}{\currentangle - \deltaangle / 2}

        % 1. 绘制扇区 (带白色描边)
        \draw[fill=\c, draw=white, thick] 
            (0,0) -- (\currentangle:\radius) arc (\currentangle:\nextangle:\radius) -- cycle;

        % 2. 处理文字旋转 (复刻 Python 里的防倒置逻辑)
        % 先将角度归一化到 [0, 360) 区间
        \pgfmathsetmacro{\normangle}{mod(\midangle + 3600, 360)}
        % 如果角度在左半边(90~270度)，则加上180度翻转，保证文字从左到右可读
        \pgfmathsetmacro{\textrot}{
            (\normangle > 90 && \normangle < 270) ? \normangle + 180 : \normangle
        }

        % 3. 写入百分比文字
        \node[font=\tiny, text=black, rotate=\textrot, anchor=center] 
            at (\midangle:\radius * \pctdistance) {\v\%};

        % 更新当前角度到下一个循环
        \global\let\currentangle\nextangle
    }

    % ======================
    % 绘制右侧图例 (复刻 bbox_to_anchor)
    % ======================
    \def\legX{4.6}    % 图例的 X 轴起点 (在半径 4 之外)
    \def\legY{3.5}    % 图例第一行的 Y 坐标
    \def\legStep{0.4} % 图例行间距

    \foreach \c/\l/\v [count=\i] in {
        c1/人与社会/18.61, c2/艺术与娱乐/12.08, c3/新闻/10.44, 
        c4/科学/8.79, c5/旅游与交通/6.43, c6/宠物与动物/5.79, 
        c7/健康/5.31, c8/敏感话题/4.91, c9/体育/4.29, 
        c10/书籍与文学/4.08, c11/法律与政府/3.86, c12/工作与教育/3.33, 
        c13/计算机与电子/2.01, c14/食品与饮料/1.96, c15/商业与工业/1.93, 
        c16/汽车与车辆/1.27, c17/游戏/1.14, c18/爱好与休闲/1.04, 
        c19/其他/2.74
    } {
        \pgfmathsetmacro{\curY}{\legY - (\i-1)*\legStep}
        % 画图例的颜色方块
        \draw[fill=\c, draw=none] (\legX, \curY) rectangle ++(0.3, 0.25);
        % 画图例的中文标签
        \node[font=\small, anchor=west] at (\legX + 0.4, \curY + 0.125) {\l};
    }

\end{tikzpicture}